\clearpage\fancyhead[L]{\small COMET.}\SourceLeaf{346}
\begin{center}\textbf{COMET.}\par\vspace{4pt}\hrule\vspace{4pt}\fontsize{10}{12}\selectfont Monatliche Correspondenz. Vol. XXIV, pp. 507\,..\,517, November 1811.\par\vspace{4pt}\hrule\vspace{4pt}\end{center}

\begin{flushright}Göttingen, 15 November.\end{flushright}
\par\indent From what \textsc{Pingré} reports of European and Chinese observations of the comet of 1301, \emph{that comet cannot be the same as the present one.} The present comet could not possibly have appeared at a longitude of 110$^\circ$ on 16 September. It is equally impossible that it appeared on 30 September at a longitude of 231$^\circ$ and a northern latitude of \emph{only} 26$^\circ$.


\par\indent Partly because of unfavorable weather and partly because of other business, I have been able to make only one observation of the comet since my last letter. I give it here.

\begin{center}
\fontsize{10}{12}\selectfont\setlength{\tabcolsep}{3pt}\renewcommand{\arraystretch}{1.15}\begin{tabular}{l|c|c}
 & R.A. & N. decl. \\ \hline
1811 Oct. 20 $10^{\mathrm u}$ $9^{\mathrm m}$ $12^{\mathrm s}$ mean time & $247^\circ\,46'\,39''$ & $41^\circ\,20'\,15''$ \\
\end{tabular}
\end{center}


\par\indent I have compared these observations with my latest parabolic elements, as Mr. \textsc{Nicolai} did with the latest observations by \textsc{Olbers} and those from Seeberg. The results follow.

\Needspace{20\baselineskip}\begin{center}Bremen observations\end{center}
\begin{center}\fontsize{10}{12}\selectfont\setlength{\tabcolsep}{8pt}\renewcommand{\arraystretch}{1.15}\begin{tabular}{l|c|c}
& \multicolumn{2}{c}{Difference} \\
Date of & & \\
observation & in R.A. & in decl. \\ \hline
Oct. 19 & $-73''$ 
 & $+152''$ \\
19 & $-73$ & ---\;--- \\
20 & $-39$ & $+139$ \\
24 & $-82$ & $+227$ \\
24 & $-84$ & $+205$ \\
25 & $-108$ & ---\;--- \\
25 & $-91$ & ---\;--- \\
28 & $-60$ & $+293$ \\
Nov. 4 & $-159$ & ---\;--- \\
4 & $-155$ & $+342$ \\
7 & $-91$ & ---\;--- \\
7 & ---\;--- & $+377$ \\
7 & $-102$ & ---\;--- \\
9 & $-274$ & $+379$ \\
\end{tabular}\end{center}
\Needspace{15\baselineskip}\begin{center}Seeberg observations\end{center}
\begin{center}\fontsize{10}{12}\selectfont\setlength{\tabcolsep}{8pt}\renewcommand{\arraystretch}{1.15}\begin{tabular}{l|c|c}
& \multicolumn{2}{c}{Difference} \\
Date of & & \\
observation & in R.A. & in decl. \\ \hline
Oct. 20 & $-113''$ & $+196''$ \\
24 & $-109$ & $+175$ \\
25 & $-88$ & $+263$ \\
28 & $-87$ & $+327$ \\
29 & $-124$ & $+265$ \\
Nov. 4 & $-110$ & $+387$ \\
5 & $-74$ & ---\;--- \\
\end{tabular}\end{center}


\Needspace{19\baselineskip}\par\indent Dr. \textsc{Olbers} has sent me Mr. \textsc{Flaugergues}’s first observations. I append Mr. \textsc{Nicolai}’s calculation of the differences from the elements

\Needspace{15\baselineskip}
\begin{center}\fontsize{9}{11}\selectfont\setlength{\tabcolsep}{3pt}\renewcommand{\arraystretch}{1.15}\begin{tabular}{l|c|c|c|c}
\shortstack{Date of\\observation} & \shortstack{Comet’s\\R.A.} & \shortstack{Deviation from\\the elements} & \shortstack{Comet’s\\S. decl.} & \shortstack{Deviation from\\the elements}\\\hline
March 26 & $120^\circ\,16'\,0''$ & $+859''$ & $29^\circ\,15'\,0''$ & $+300''$\\
28 & $119\;52\;56$ & $+316$ & $28\;\;7\;\;0$ & $+257$\\
29 & $119\;41\;\;4$ & $+127$ & $27\;32\;57$ & $+226$\\
30 & $119\;29\;26$ & $-50$ & $26\;58\;22$ & $+193$\\
31 & $119\;18\;36$ & $-306$ & $26\;23\;13$ & $+297$\\
April 1 & $119\;\;7\;38$ & $-382$ & $25\;50\;17$ & $+219$\\
\end{tabular}\end{center}

\SourceLeaf{347}\par\indent In earlier reports the right ascension of Mr. \textsc{Flaugergues}’s first observation was given as $10'$ larger, in which case the difference would be only $259''$. I am withholding judgment on the large discrepancies until I find time to do calculations on the comet again myself, something I still cannot get to just now.


